\section*{第 \ref{ch:g10:light-spectra} 章 --- 光谱与光的信息}

\begin{solution}{exo:g10:light-spectra:1}
$\qty{254}{nm} < \qty{400}{nm}$：紫外线。$\qty{480}{nm}$：可见光，
蓝色。$\qty{589}{nm}$：可见光，黄橙色。$\qty{656}{nm}$：可见光，
红色。$\qty{1064}{nm} > \qty{800}{nm}$：红外线 --- 这使切割用激光器
具有双重危险：功率强大，且不可见。
\end{solution}

\begin{solution}{exo:g10:light-spectra:2}
\emph{1.} 由红到紫的连续色带：
太阳光的光谱。

\emph{2.} 紫端 --- 棱镜对短波长
偏折最大。

\emph{3.} 单独一个绿色光点：发生了偏折，但未展开成光带。
激光是单色的 --- 只有一个波长，故无
需要分开的成分。
\end{solution}

\begin{solution}{exo:g10:light-spectra:3}
(a) 连续光谱（致密的热灯丝）。(b) 发射光谱（低压
受激气体）。(c) 连续光谱（致密的熔融金属）。(d) 发射光谱
（实际上仅有 $\qty{589}{nm}$ 处的一条谱线）。(e) 吸收光谱：热表面的
连续光谱上，有太阳较冷大气层的
暗线。
\end{solution}

\begin{solution}{exo:g10:light-spectra:4}
\emph{1.} 雨滴：每一滴接收白色阳光，将其分解
并按颜色分开后射回。

\emph{2.} 太阳提供白光；雨水提供数以百万计的
微小棱镜；几何关系只对由雨滴射向观察者的
光成立，这也是为什么太阳必须在观察者身后，
雨在前方。

\emph{3.} 雨滴本身不发光 --- 它们只是分开所接收到的
光。若分开后的阳光呈现出所有颜色，则这些颜色
原本就包含于阳光之中：阳光是白光。
\end{solution}

\begin{solution}{exo:g10:light-spectra:5}
\emph{1.} 暗红 $<$ 亮橙 $<$ 近白。

\emph{2.} 两者兼有：发光\emph{更亮}（效应 1），其颜色
向蓝端漂移（效应 2，峰值 $\lambda_{\max}$
减小），因此可见光带中越来越多的部分被强烈照亮 ---
整个光带便呈白色。

\emph{3.} 发白光的物体在整个可见光
波段内强烈发射；发红光的物体仅在长波端有发射。由维恩
定律可知，前者温度更高。
\end{solution}

\begin{solution}{exo:g10:light-spectra:6}
\emph{1.}
\[
T = \frac{\qty{2.90e-3}{m.K}}{\qty{5.00e-7}{m}} = \qty{5800}{K}.
\]

\emph{2.} $\theta = 5800 - 273 = \qty{5527}{\degreeCelsius} \approx
\qty{5500}{\degreeCelsius}$.

\emph{3.} $\qty{400}{nm} < \qty{500}{nm} < \qty{800}{nm}$：峰值
位于可见光的绿色区域内。
\end{solution}

\begin{solution}{exo:g10:light-spectra:7}
\emph{1.} 蜡烛：
$\lambda_{\max} = \qty{2.90e-3}{m.K} / \qty{1800}{K} = \qty{1.61e-6}{m}
= \qty{1.6}{\micro m}$，位于红外区。电热板：
$\lambda_{\max} = \qty{2.90e-3}{m.K} / \qty{500}{K} = \qty{5.8e-6}{m}
= \qty{5.8}{\micro m}$，远红外区。

\emph{2.} 在 $\qty{500}{K}$ 时，辐射几乎全为红外线：
皮肤吸收它并感到热，但眼睛接收不到任何光 ---
热却是黑的。

\emph{3.} 连续光谱很宽：即使峰值位于
$\qty{1.6}{\micro m}$ 处，蜡烛光谱在
红黄端仍有可见的尾部 --- 虽然微弱，但对暗适应的眼睛已足够。
\end{solution}

\begin{solution}{exo:g10:light-spectra:8}
灯 A：谱线 $434$、$486$、$\qty{656}{nm}$ 与氢吻合，另有
暗弱的紫色谱线在 $\qty{410}{nm}$ 处补全谱线表：氢。

灯 B：谱线位于 $\qty{589}{nm}$，是钠唯一的强可见
谱线。氦并非因这条谱线本身被排除（氦在邻近的
$\qty{588}{nm}$ 处确有一条谱线），而是因\emph{缺失}的谱线被排除：氦还应出现
$447$、$502$ 和 $\qty{668}{nm}$，而这些均未出现。
元素由完整谱线表确定：灯 B 含有钠。
\end{solution}

\begin{solution}{exo:g10:light-spectra:9}
\emph{1.} 一条连续的彩虹色带，上有一条细暗线：
钠的吸收光谱。

\emph{2.} 在 $\qty{589}{nm}$ 处：气体恰好吸收它能够
发射的波长（见\cref{prop:g10:light-spectra:kirchhoff}），钠的
强可见谱线位于 $\qty{589}{nm}$。

\emph{3.} 钠的发射光谱：一条明亮的黄橙色
谱线，置于黑色背景上，位于 $\qty{589}{nm}$ ---
波长相同，两种光谱的明暗恰好互换。
\end{solution}

\begin{solution}{exo:g10:light-spectra:10}
\emph{1.} $410$、$434$、$486$、$\qty{656}{nm}$：完整的氢
谱线表。$517$ 和 $\qty{518}{nm}$：镁的双线。两种元素均
存在。

\emph{2.} 太阳大气层中，位于发光表面上方的
较冷、较稀薄的气体。表面发出连续光谱；
大气层从穿过的光中吸收自身的波长，因此谱线
表现为缺失的光 --- 在明亮背景上呈暗线。
\end{solution}

\begin{solution}{exo:g10:light-spectra:11}
\emph{1.} 取 $T = 37 + 273 = \qty{310}{K}$，则
\[
\lambda_{\max} = \frac{\qty{2.90e-3}{m.K}}{\qty{310}{K}}
= \qty{9.4e-6}{m} = \qty{9.4}{\micro m}.
\]

\emph{2.} 远红外 ---
波长是红光波长的十倍以上。

\emph{3.} 在 $\qty{310}{K}$ 时，辐射中的可见光部分
完全可以忽略：无论房间多暗，都没有可见的辉光。
热成像仪的传感器调谐于实际发出的、
约 $\qty{9}{\micro m}$ 的辐射，
并能看到较冷背景中发光的温暖物体。
\end{solution}

\begin{solution}{exo:g10:light-spectra:12}
\emph{1.} $\lambda_{\max} = \qty{2.90e-3}{m.K} / \qty{2700}{K}
= \qty{1.07e-6}{m} \approx \qty{1.1}{\micro m}$：近红外。

\emph{2.} 光谱峰值位于可见光范围之外：大部分电功率
以看不见的红外线形式辐射出去，即热量，只有
很小一部分可见光尾部照亮房间。作为灯具，灯泡浪费了
大部分电功率；作为加热器，它几乎是完美的。

\emph{3.} $T = \qty{2.90e-3}{m.K} / \qty{5.50e-7}{m} \approx
\qty{5300}{K}$ --- 远高于钨的熔点，
$\qty{3700}{K}$。无任何固体灯丝能被加热到日光般的温度，
因此日光型灯具不得不放弃白炽发光，
（荧光灯管、发光二极管）产生可见光，
而无需把物体加热到 $\qty{5300}{K}$。
\end{solution}

\begin{solution}{exo:g10:light-spectra:13}
\emph{1.} $T = \qty{2.90e-3}{m.K} / \qty{4.20e-7}{m} \approx
\qty{6900}{K}$：比太阳的 $\qty{5800}{K}$ 更热。

\emph{2.} 氢（完整的一组 $410$、$434$、$486$、$\qty{656}{nm}$ 均
存在）和钙（两条线，$393$ 与 $\qty{397}{nm}$）。

\emph{3.} 否。$\qty{589}{nm}$ 处的钠谱线不存在，若大气中
含有大量的钠，必定会在连续背景上留下该谱线。
若确有钠存在，也仅是痕量，
微弱到无法探测。
\end{solution}

\begin{solution}{exo:g10:light-spectra:14}
\emph{1.} 氢可解释 $410$、$434$、$486$ 和 $\qty{656}{nm}$；
氦可解释 $447$、$502$、$588$ 和 $\qty{668}{nm}$。所观测到的每条谱线
均有所归属：该混合物为氢和氦。

\emph{2.} 氦可以作为\emph{充分}的解释，因为它的其余三条
谱线（$447$、$502$、$\qty{668}{nm}$）均存在 ---
$\qty{588}{nm}$ 谱线此时仅凭氦即可预期，光谱中
并无任何内容\emph{要求}存在钠。要定论此事，所需的是
分辨率而非争辩：更精细的分光镜能将氦的单线
（位于 $\qty{587.6}{nm}$）与钠的谱线区分开来，钠的谱线实为一对
位于 $589.0$ 与 $\qty{589.6}{nm}$ 的靠近谱线。此处只见一条线则可排除钠；
见三条则可定钠的罪。
\end{solution}

\begin{solution}{exo:g10:light-spectra:15}
\emph{1.} 参宿四：
$T = \qty{2.90e-3}{m.K} / \qty{8.5e-7}{m} \approx \qty{3400}{K}$。参宿七：
$T = \qty{2.90e-3}{m.K} / \qty{2.4e-7}{m} \approx \qty{12000}{K}$。

\emph{2.} 由于两者的 $\lambda_{\max} \times T$ 为同一常数，
$\dfrac{T_{\text{参宿七}}}{T_{\text{参宿四}}}
= \dfrac{\qty{850}{nm}}{\qty{240}{nm}} \approx 3.5$。

\emph{3.} 参宿四的峰值位于红外：在可见光波段内，其
光谱在红端最强，因此这颗星看起来呈红色。参宿七的峰值位于
紫外：其可见光部分在蓝端最强，因此呈
蓝白色。

\emph{4.} 连续光谱覆盖极宽的波长范围。即使
峰值位于可见光波段之外，恒星仍会
向该波段大量\emph{射入}光 --- 决定色调的是峰值，而非能否看见。
\end{solution}

\begin{solution}{pb:g10:light-spectra:1}
\textbf{1.} 约 $\qty{400}{nm}$（紫端）至
约 $\qty{800}{nm}$（红端）；$\qty{550}{nm}$ 位于中部，属于
绿色。

\textbf{2.} $\qty{410}{nm}$：紫色；$\qty{434}{nm}$：蓝紫色；
$\qty{486}{nm}$：蓝绿色；$\qty{656}{nm}$：红色。

\textbf{3.} 单色的黄橙色光辉。汽车的颜色取决于
车漆所反射的各种波长的混合；在只提供
\emph{仅}$\qty{589}{nm}$ 光的灯下，每种漆都只能或多或少地反射
这一种波长 --- 于是所有汽车都只呈现橙色和
灰色的深浅。

\textbf{4.} 每种元素发射固定的谱线组，
任何两种元素的谱线组都不相同
（见\cref{prop:g10:light-spectra:fingerprint}）：因此一组完整的谱线
可像指纹一样标识元素。但
\emph{部分}匹配说明不了什么 --- 不同元素可能拥有几乎重合的单线
（氦位于 $\qty{588}{nm}$，钠位于
$\qty{589}{nm}$）--- 因此断定某元素存在，须要求
谱线表中的所有强线均出现。

\textbf{5.} (a) 连续彩虹带，在
$\qty{589}{nm}$ 处有一条暗线：吸收光谱。(b) 在
黑底上、$\qty{589}{nm}$ 处有一条亮线：发射光谱 --- 波长相同，
明暗恰好互换
（见\cref{prop:g10:light-spectra:kirchhoff}）。

\textbf{6.} 光谱在每一波长处都增强，其峰值
$\lambda_{\max}$ 向较短波长方向移动，遵循
$\lambda_{\max} \times T = \qty{2.90e-3}{m.K}$
（见\cref{prop:g10:light-spectra:wien}）。

\textbf{7.} $\lambda_{\max} = \qty{2.90e-3}{m.K} / \qty{3400}{K}
= \qty{8.5e-7}{m} = \qty{850}{nm}$：近红外，刚越过红光
边缘 --- 这正是卤素灯色调比日光偏暖的原因。

\textbf{8.} $T = \qty{2.90e-3}{m.K} / \qty{5.00e-7}{m} = \qty{5800}{K}$.

\textbf{9.} 峰值平缓而非尖锐：太阳在整个
可见光波段上发射的强度相当。各色光混合在一起
呈白色 --- 太阳看上去是白的（经大气层后略带黄色），
绝不会是绿色。

\textbf{10.} 冷火钳：峰值位于深红外区，没有可见光尾部 ---
可感受到辐射，却看不到光。较热时：光谱增强，
出现第一段可见光尾，超过 $\qty{800}{nm}$ --- 红端
单独发亮：暗红色。更热时：整个可见光波段逐渐填满，
同时整体变亮：先呈橙色，再趋向白色。

\textbf{11.} 连续背景来自致密、炽热的发光
表面；暗线则由其上方较冷、较稀薄的大气层造成，
它从穿过的光线中吸收自身的波长，
故呈暗线。

\textbf{12.} 氢（$410$、$434$、$486$、$\qty{656}{nm}$：完整）、
钠（$\qty{589}{nm}$）和钙（$393$ 与 $\qty{397}{nm}$：两者皆有）
--- 三种元素均存在于太阳大气中。

\textbf{13.} 大气层以氢为主。现代测量与此相符，并给出定量结果：
按质量计，太阳约有四分之三是氢，
其余大部分为氦。

\textbf{14.} 该谱线与任何已知元素都不对应；由于每种元素的谱线表固定且唯一，
不属于任何谱线表的谱线，必定属于实验室中从未见过的某种元素。
于是宣布发现了一种新元素，
并以希腊语 \emph{helios}（太阳）命名为氦。
它于 1895 年在地球上被分离 --- 该预言
在二十七年后得到证实。

\textbf{15.} 在明亮的连续背景下，
大气层在其自身的波长处\emph{吸收}光：呈现暗线。
全食期间月球遮挡了背景；稀薄的大气层随即
独自映衬在暗天空中，其再发出的光便是唯一可见之物：
波长相同，此时却呈亮线。同一气体、同一谱线表
--- 明暗取决于其后方是什么。

\textbf{16.} $T = \qty{2.90e-3}{m.K} / \qty{2.90e-7}{m}
= \qty{10000}{K}$.

\textbf{17.} 蓝白色：在可见光波段内，光谱在蓝端最强。
也并非不可见：连续光谱充满整个可见光波段，
途经可见光后抵达紫外峰值 ---
峰值决定色调，而非亮度。

\textbf{18.} 暗线 $410$、$434$、$486$、$\qty{656}{nm}$ 构成
完整的氢指纹：织女星大气以氢为主。

\textbf{19.} $T_{\text{织女星}} / T_{\text{太阳}} = 10000 / 5800 \approx
1.7$：织女星温度更高，因而更偏蓝（峰值位于 $\qty{290}{nm}$，
而太阳为 $\qty{500}{nm}$）。由
\cref{prop:g10:light-spectra:wien}的第一部分可知，单位发光表面在
每一波长处的亮度也更高。

\textbf{20.} 织女星身份卡 --- 表面温度约
$\qty{10000}{K}$；大气层的主导元素：氢；视觉颜色：
蓝白色。以上全部信息都由一束光读出，借助棱镜、
实验室谱线表与一次除法。孔德的“永远不可能”
提出于 1835 年，1859 年失效：持续了二十四年。
\end{solution}
